\documentclass{exam}
\usepackage[utf8]{inputenc}
\usepackage[spanish]{babel}
\usepackage{amsmath}


\pointpoints{punto}{puntos}

\newcommand{\m}[5]{
  \begin{tabular}{p{3cm} p{3cm} p{3cm} p{3cm} p{3cm}}
a) #1 &b) #2 &c) #3 &d) #4 &e) #5 \\ 
  \end{tabular}
}

\begin{document}
\begin{center}
    \textbf{\large Evaluación de Electromagnetismo}
\end{center}

\begin{questions}

\question[1] El vacío tiene una permeabilidad magnética $\mu_0 = 4\pi \times 10^{-7}$ T$\cdot$m/A. ¿Cuál es la permeabilidad magnética relativa $\mu_r$ de un material que tiene una permeabilidad magnética de $1.8 \times 10^{-5}$ T$\cdot$m/A?

\m{2.25}{4}{6.25}{1.43}{0}


\question[1] Un solenoide de 200 vueltas y longitud de 30 cm tiene una corriente eléctrica de 2 A. ¿Cuál es la inducción magnética en el centro del solenoide?

\m {0.002 T}{0.04 T}{0.2T}{2T}{4.23T}

 \question[2] La intensidad de campo magnético en el interior del solenoide es de $H=200\frac{Av}{m}$ y su longitud es de $l=30cm$. Calcular:
    \begin{enumerate}
        \item Amperio vuelta del solenoide.
        
        \m{40A}{60A}{80A}{100A}{120A}
        
        \item Intensidad de corriente $I$, en el conductor si el solenoide  tiene $N=1000$ espiras.
        
        \m{0,167A}{0,6A }{0,06A}{0.7A}{1 A}
    \end{enumerate}

\end{questions}

\end{document}
